\subsection{The continuous unit-equation obstruction}

\begin{statement}{A.25}{Proved}{finite-coefficient vanishing for the actual ghost locus}
Let $\mathcal G$ be the set of pairs $(A,B)\in
\widehat{\Q^\times}^{\,2}$ such that, for each finite prime $p$, there is
$x_p\in\Q_p\setminus\{0,1\}$ with
$A_p=\kappa(x_p)$ and $B_p=\kappa(1-x_p)$.
For every integer $n\geq1$ the image of $A\cup B$ in
$H^2(G_\Q,\mu_n^{\otimes2})$ is zero. The continuous class lies in the
kernel in the exact sequence
\[
0\longrightarrow\varprojlim_m^{\,1}
 H^1(G_\Q,\mu_{m!}^{\otimes2})
 \longrightarrow H^2_{\rm cts}(G_\Q,\widehat{\Z}(2))
 \longrightarrow\varprojlim_m H^2(G_\Q,\mu_{m!}^{\otimes2})
 \longrightarrow0.
\]
Its restrictions to all continuous local cohomology groups, including
the real one, are zero. No local section theorem is used.
\end{statement}
\begin{proof}
We first prove the global detection assertion needed here. The exact
tame-symbol sequence and the integral calculation are
\[
0\longrightarrow K_2(\Z)\longrightarrow K_2(\Q)
 \xrightarrow{\oplus\partial_p}\bigoplus_p\mathbf F_p^\times
 \longrightarrow0,
\qquad K_2(\Z)=\langle\{-1,-1\}\rangle\simeq\Z/2.
\]
These are the specialization to $\Z$ of the localization theorem and
the calculation of $K_2(\Z)$ in
\cite[III.5.2.2, III.6.5--6.5.1]{Weibel2013}.
The Galois symbol identifies $K_2(F)/n$ with
$H^2(G_F,\mu_n^{\otimes2})$ for each field of characteristic zero;
see \cite[III.6.10.3--6.10.4, VI.4.1.1]{Weibel2013}.

Let $z\in K_2(\Q)$ have image zero in every $K_2(\Q_p)/n$.
The finitely supported vector $(\partial_p z)_p$ belongs to
$n\bigoplus_p\mathbf F_p^\times$: each coordinate does, by functoriality
of the tame symbol. Choose a finitely supported $n$th-root vector and
lift it, by surjectivity, to $w\in K_2(\Q)$. Then $z-nw$ belongs to
$K_2(\Z)$. For odd $n$ this group contributes zero modulo $n$.
For even $n$ its nonzero element is detected already by the quadratic
Hilbert symbol over $\Q_2$. Indeed the norm equation
$X^2+Y^2=-1$ has no $\Q_2$-solution: on clearing a common denominator
and reducing a primitive triple modulo $4$, a sum of three squares
cannot be zero. Thus $\{-1,-1\}$ is not divisible by $2$ in
$K_2(\Q_2)$, and cannot be divisible by even $n$. This proves that
\[
H^2(G_\Q,\mu_n^{\otimes2})
 \longrightarrow\prod_{p<\infty}H^2(G_{\Q_p},\mu_n^{\otimes2})
\]
is injective.

For $(A,B)\in\mathcal G$, localization of its finite-level cup is
$\kappa_n(x_p)\cup\kappa_n(1-x_p)=0$ by the Steinberg relation.
The just-proved injectivity gives global vanishing. The coefficient
transition maps $\mu_{(m+1)!}^{\otimes2}\to\mu_{m!}^{\otimes2}$
are the standard reductions for Tate twists. Their continuous cochain
complexes form a degreewise surjective tower: a continuous map into a
finite discrete coefficient group lifts by choosing a lift for each
of its finitely many values. Applying the inverse-limit exact sequence
to these complexes gives the displayed sequence. This is also
\cite[(2.1), Theorem (2.2)]{Jannsen1988} with the cofinal factorial
tower. For every local field and every finite coefficient module the
$H^1$ groups are finite. Their inverse systems therefore satisfy the
Mittag--Leffler condition, so the corresponding local $\varprojlim^1$
is zero. Finite-level vanishing proves continuous local vanishing.
The same argument applies to the order-two real Galois group.
\end{proof}

\begin{statement}{A.26}{Proved}{the alternating continuous residue at the coefficient prime three}
Write
\[
R_3=\{(z_q)_{q\ne3}\in\prod_{q\ne3}\Z_3:
 z_q\longrightarrow0\text{ in the cofinite sense}\},
\qquad d(z)_q=(q-1)z_q.
\]
For $(A,B)\in\mathcal G$, choose local tame Frobenius lifts at $q\ne3$
and integral Kummer cocycles. Let their Frobenius values be
$\alpha_q,\beta_q\in\Z_3$, and let their inertia values be the ordinary
integer valuations $a_q=v_q(x_q)$ and $b_q=v_q(1-x_q)$.
Then the vector
\[
 z_q=\alpha_q b_q-\beta_q a_q
\]
belongs to $dR_3$. In particular its class in $R_3/dR_3$ is zero.
\end{statement}
\begin{proof}
The valuation vectors are the $3$-components of the global Kummer
valuation vectors, so both tend to zero $3$-adically. At each $q$ the
ordinary local equation has exactly four possible valuation patterns.
For $a_q>0$ we have $b_q=0$ and $1-x_q\in1+q\Z_q$.
This pro-$q$ unit has zero pro-$3$ Kummer class, so
$\beta_q\in(q-1)\Z_3$. Hence
$z_q/(q-1)=-a_q\beta_q/(q-1)\in a_q\Z_3$.
For $b_q>0$ the symmetric argument gives
$z_q/(q-1)\in b_q\Z_3$.
For $a_q<0$, the equality $b_q=a_q$ holds and
$(1-x_q)/x_q\in-1\cdot(1+q\Z_q)$.
The pro-$3$ Kummer class of $-1$ is zero, as is that of the indicated
principal unit. Thus
$\beta_q-\alpha_q\in(q-1)\Z_3$ and
$z_q/(q-1)\in a_q\Z_3$.
Finally $a_q=b_q=0$ gives $z_q=0$.
The quotient vector belongs to $R_3$ in every case because $a_q,b_q$
tend to zero. Multiplication by $q-1$ is injective on $\Z_3$, so these
quotients are unambiguous. Changing a local Kummer cocycle by a
coboundary changes its Frobenius value by an element of
$(q-1)\Z_3$; the resulting quotient difference is again a multiple of
$a_q$ or $b_q$, hence belongs to $R_3$.
\end{proof}

% Replaces the former open A.27 and its following computation.
% Uses the ordinary local unit-equation locus \mathcal G of A.25.
\begin{statement}{A.27}{Proved}{The full continuous cup vanishes on the ordinary local unit-equation locus}
For every $(A,B)\in\mathcal G$,
\[
 A\cup B=0
 \quad\text{in }H^2_{\mathrm{cts}}(G_\Q,\widehat\Z(2)).
 \tag{A.27.1}
\]
This includes the $2$-primary component and all first derived
inverse-limit components. The assertion imposes no finite-support
condition on the valuation vectors. It constructs no
infinite-support element of $\mathcal G$ and makes no assertion
about the obstruction in weight three.
\end{statement}
\begin{proof}
Recall that $\mathcal G$ consists of pairs in
$\widehat{\Q^\times}^{\,2}$ for which, at every finite prime $q$,
there is an ordinary $x_q\in\Q_q\setminus\{0,1\}$ with
$A_q=\kappa(x_q)$ and $B_q=\kappa(1-x_q)$.
We compute the actual continuous residue class using a strict
map of cochain complexes. Finite-level vanishing alone would
leave the first derived inverse limit uncontrolled.

Fix a prime $\ell$. Put $T=\Z_\ell(1)$,
$M=\Z_\ell(2)$, $M_n=M/\ell^nM$, and $\chi=\chi_\ell$.
Choose compatible bases of the Tate modules and normalized
continuous cocycles $a,b:G_\Q\to T$ representing $A_\ell,B_\ell$.
The convention for differentials and the cup representative is
\[
 (\delta f)(g,h)=g f(h)-f(gh)+f(g),\qquad
 c(g,h)=a(g)\chi(g)b(h).
 \tag{A.27.2}
\]
We use this original cup cocycle integrally at every prime.

\emph{The exact global localization complex.}
Write $S_\ell=\operatorname{Spec}\Z[1/\ell]$ and
$G_{\Q,S_\ell}$ for the Galois group unramified at all finite
primes different from $\ell$. The arithmetic--\'etale comparison
is \cite[II, Proposition 2.9]{BIR_MilneADT}; the finite
cohomology groups for finite coefficients are
\cite[I, Corollary 4.15]{BIR_MilneADT}.
Localization over a finite set of closed points of $S_\ell$,
followed by its filtered union, gives a coefficient-compatible
triangle
\[
 R\Gamma(S_\ell,M_n)\longrightarrow R\Gamma(\Q,M_n)
 \longrightarrow
 \bigoplus_{q\ne\ell}R\Gamma(\mathbf F_q,M_n(-1))[-1]
 \longrightarrow R\Gamma(S_\ell,M_n)[1].
 \tag{A.27.3}
\]
The local term is the codimension-one purity term; its twist
and support maps are the case used in
\cite[\S3, (3.15)--(3.20)]{BIR_Jannsen1988}.
Each residue field has procyclic Galois group, and its
Frobenius acts on $M_n(-1)$ by $q$. Thus the unshifted third
term has the compatible two-term representative
\[
 C_{\ell,n}=
 \left[\bigoplus_{q\ne\ell}\Z/\ell^n
       \xrightarrow{\ d\ }
       \bigoplus_{q\ne\ell}\Z/\ell^n\right],
 \qquad d(v)_q=(q-1)v_q,
 \tag{A.27.4}
\]
in degrees zero and one.

Both termwise transition maps are onto. Define the literal
inverse-limit module
\[
 R_\ell=
 \{(v_q)_{q\ne\ell}\in\prod_{q\ne\ell}\Z_\ell:
       v_q\longrightarrow0\text{ in the cofinite sense}\}.
 \tag{A.27.5}
\]
It is exactly $\varprojlim_n\bigoplus_{q\ne\ell}\Z/\ell^n$:
finite support modulo each $\ell^n$ is the displayed
convergence condition. Therefore
$R\varprojlim_n C_{\ell,n}=[R_\ell\xrightarrow d R_\ell]$.
The continuous global cochain complexes are also the derived
inverse limits of their finite-coefficient complexes. Indeed
each finite-valued continuous cochain lifts continuously
through a surjection of coefficient groups by lifting its
finitely many values. This is the continuous-cochain comparison
in \cite[(1.6), (2.1), Theorem (2.2)]{BIR_Jannsen1988}.
Taking $R\varprojlim$ in (A.27.3) yields, in particular, the
exact sequence
\[
 H^2_{\mathrm{cts}}(G_{\Q,S_\ell},M)
 \longrightarrow H^2_{\mathrm{cts}}(G_\Q,M)
 \xrightarrow{\ \mathrm{res}\ } R_\ell/dR_\ell.
 \tag{A.27.6}
\]
No vanishing of a real higher-cohomology group is needed for
this three-term exact sequence. Its denominator is the actual
image $dR_\ell$, without taking its closure.

\emph{A strict residue map, including its normalization.}
At $q\ne\ell$, discard wild inertia and the prime-to-$\ell$
part of tame inertia. Their positive-degree $\ell$-primary
cohomology vanishes. Here is a strict cochain map implementing
this passage. Write $N$ for their combined kernel and use
normalized Haar measure on the pro-prime-to-$\ell$ group $N$.
For a homogeneous continuous cochain $f$ with values in $M_n$,
set
\[
 (P_nf)(\bar g_0,\ldots,\bar g_r)
   =\int_{N^{r+1}} f(g_0n_0,\ldots,g_rn_r)\,
                                  dn_0\cdots dn_r.
 \tag{A.27.7a}
\]
Each integral is a finite average whose denominator is prime
to $\ell$. It is independent of the lifts $g_i$, is equivariant
for the quotient group, and commutes with the homogeneous
cochain differential, since that differential alternately
omits one argument and the omitted Haar factor has integral
one. The maps $P_n$ commute with coefficient reduction and
are the identity on inflated cochains. Inflation is a
quasi-isomorphism by the vanishing of the positive-degree
cohomology of $N$; hence $P_n$ is its strict cochain inverse
up to quasi-isomorphism. The homogeneous-to-inhomogeneous
identification sends $f$ to
$u(g_1,\ldots,g_r)=f(1,g_1,g_1g_2,\ldots,g_1\cdots g_r)$.
If averaging has destroyed normalization, use in the two
degrees required below the explicit projections
\[
 \begin{split}
 (N^1u)(g)&=u(g)-u(1),\\
 (N^2v)(g,h)&=v(g,h)-v(g,1)-g v(1,h)+g v(1,1).
 \end{split}
 \tag{A.27.7b}
\]
Their values vanish when an argument is $1$, and direct
substitution gives $N^2\delta=\delta N^1$ and
$N^1\delta=\delta$ on zero-cochains. They fix normalized
cochains. Thus their composition with $P_n$ is a
coefficient-compatible strict map in exactly the degrees
used by the residue complex, and it fixes all inflated
normalized unramified cochains. We may compute on
\[
 D_q=\Z_\ell(1)\rtimes\widehat\Z,
 \qquad F\tau F^{-1}=\tau^q.
 \tag{A.27.7}
\]
Here $F$ is arithmetic Frobenius, and $\tau$ is normalized by
the Kummer character of the uniformizer $q$. It acts trivially
on $M$, while $F$ acts by $q^2$. The cocycles $a,b$ and $c$
already factor through (A.27.7) on the local group, since
their restrictions to the discarded inertia vanish.

For normalized local cochains of degrees one and two, set
\[
 \begin{split}
 \psi_q^1(u)&=u(\tau),\\
 \psi_q^2(v)&=q^{-1}\left(
 v(F,\tau)+v(F\tau,F^{-1})-v(F,F^{-1})
       -\sum_{j=1}^{q-1}v(\tau^j,\tau)\right).
 \end{split}
 \tag{A.27.8}
\]
These formulas are integral because $q$ is an $\ell$-adic
unit. All choices are made once before reduction modulo
$\ell^n$. On $v=\delta u$, the first three terms telescope
to $q^2u(\tau)-u(\tau^q)$, and the sum telescopes to
$q u(\tau)-u(\tau^q)$. Hence
\[
 \psi_q^2(\delta u)=(q-1)\psi_q^1(u).
 \tag{A.27.9}
\]
Also $\psi_q^1(\delta m)=0$ for a zero-cochain $m$.
Taking the other components to be zero gives a strict map to
$[M\xrightarrow{q-1}M]$ in degrees one and two.

Both formulas vanish on normalized cochains inflated from
$\widehat\Z$. They therefore descend to the literal quotient
of the local cochain complex by its unramified subcomplex.
The latter represents the local localization cone. On
degree-one cohomology the map is evaluation of an inertia
character. On degree-two cohomology it evaluates the error
in lifting $F\tau F^{-1}=\tau^q$, divided by $q$.
The two-row Hochschild--Serre calculation, since both
$\Z_\ell(1)$ and $\widehat\Z$ have $\ell$-cohomological
dimension one, identifies the two relative groups with
the kernel and cokernel of $q-1$. Thus (A.27.8) is a
quasi-isomorphism to the purity term in (A.27.3).
The factor $q^{-1}$ is the Tate twist: Frobenius transports
a character of inertia by evaluating it on
$\tau^{q^{-1}}$, and hence acts on
$\operatorname{Hom}(\Z_\ell(1),M)$ by $q^2/q=q$.

This also fixes the global chain comparison. For a finite
ramification set $J$, restriction maps the cone of the
unramified global cochains into the sum, over $q\in J$,
of these relative local quotients, followed by the maps
$P_n$ of (A.27.7a), the projections (A.27.7b), and (A.27.8).
The unramified component
is killed identically by the quotient and by (A.27.8).
Excision identifies this map with (A.27.3). Each normalized
finite-coefficient global cochain factors through a finite
Galois quotient, so outside a finite $J$ its local restriction
is inflated from the unramified quotient. Its image under
(A.27.8) is then literally zero. The construction therefore
lands in the direct sums at finite level and commutes
strictly with every coefficient transition. Its inverse
limit determines the actual continuous residue class,
including its first derived inverse-limit component.

\emph{The residue of the original cup.}
Write
\[
 a(\tau)=a_q,\quad b(\tau)=b_q,\quad
 a(F)=\alpha_q,\quad b(F)=\beta_q.
\]
The local Kummer identifications give
$a_q=v_q(x_q)$ and $b_q=v_q(1-x_q)$ as ordinary integers.
For every fixed coefficient level the corresponding inertia
characters vanish outside a finite set; consequently
$a_q,b_q\to0$ in $\Z_\ell$ as $q$ leaves finite sets.
The cocycle identities give
\[
 \begin{gathered}
 a(F\tau)=\alpha_q+q a_q,
 \qquad b(F^{-1})=-q^{-1}\beta_q,\\
 c(F,\tau)=q\alpha_q b_q,
 \quad c(F\tau,F^{-1})=-\alpha_q\beta_q-q a_q\beta_q,
 \quad c(F,F^{-1})=-\alpha_q\beta_q,\\
 c(\tau^j,\tau)=j a_qb_q.
 \end{gathered}
\]
Substitution in (A.27.8) proves the exact formula
\[
 z_q:=\psi_q^2(c)
   =\alpha_q b_q-\beta_q a_q-
                          \frac{q-1}{2}a_qb_q,
 \qquad \mathrm{res}([c])=[(z_q)]\in R_\ell/dR_\ell.
 \tag{A.27.10}
\]
For odd $\ell$, the displayed half is integral in $\Z_\ell$.
For $\ell=2$, $q$ is odd, so $(q-1)/2$ is an ordinary
integer. The finite sum in (A.27.8) proves this formula
without dividing a $2$-adic cochain by two.

For comparison with A.26, at an odd coefficient prime the
continuous one-cochain $f(g)=a(g)b(g)/2$ satisfies
\[
 c+\delta f
   =\tfrac12\chi(g)\bigl(a(g)b(h)-b(g)a(h)\bigr).
\]
Its residue is $\alpha_qb_q-\beta_qa_q$, because
$\psi_q^2(\delta f)=(q-1)a_qb_q/2$.
The family $(a_qb_q/2)_q$ belongs to $R_\ell$.
Thus A.26 computes precisely the alternating representative
at $\ell=3$; (A.27.10) gives the original integral cup at
every coefficient prime.

\emph{A primitive that lies in the restricted module.}
There are four valuation patterns for the ordinary local
equation $x_q+(1-x_q)=1$.

If $a_q>0$, then $b_q=0$ and $1-x_q\in1+q\Z_q$.
This pro-$q$ unit has zero $\ell$-primary Kummer class for
$q\ne\ell$. Thus $\beta_q=(q-1)u_q$ for an integral
$u_q\in\Z_\ell$, and
$z_q/(q-1)=-a_qu_q$.
If $b_q>0$, interchange the two roles: $a_q=0$,
$\alpha_q=(q-1)u_q$, and $z_q/(q-1)=b_qu_q$.
If $a_q=b_q=0$, then $z_q=0$.

In the remaining case $a_q=b_q=m_q<0$ and
\[
 \frac{x_q}{1-x_q}=-\left(1-\frac1{x_q}\right)^{-1}.
 \tag{A.27.11}
\]
The factor in parentheses is a principal unit. If $\ell$ is
odd, $-1$ also has trivial $\ell$-primary Kummer class, so
$\alpha_q-\beta_q=(q-1)u_q$ with $u_q\in\Z_\ell$ and
\[
 \frac{z_q}{q-1}=m_qu_q-\frac{m_q^2}{2}.
 \tag{A.27.12}
\]
For $\ell=2$, a compatible choice of $2^n$th roots of
$-1$ is $\zeta_{2^{n+1}}$, with the compatible basis
$\zeta_{2^n}$ for the Tate module. Its actual Kummer
cocycle is
\[
             \kappa_2(-1)(g)=\frac{\chi_2(g)-1}{2}.
 \tag{A.27.13}
\]
This is integral, since every $\chi_2(g)$ is odd.
Equation (A.27.11) now gives, with the sign fixed by
(A.27.13),
\[
 \alpha_q-\beta_q=(q-1)\left(u_q+\tfrac12\right),
 \qquad
 \frac{z_q}{q-1}=m_qu_q-\frac{m_q(m_q-1)}2
                =m_qu_q-\binom{m_q}{2}.
 \tag{A.27.14}
\]
The binomial expression is integral for every integer $m_q$,
including negative integers, and indeed for every $m_q\in\Z_2$.
It has $2$-adic valuation at least $v_2(m_q)-1$.

Define $w_q=z_q/(q-1)$ using these formulas. All its
coordinates are integral. In the first two patterns they
are an integral multiple of $a_q$ or $b_q$. In the
negative pattern (A.27.12) gives the same bound at odd
$\ell$, while (A.27.14) loses at most one power of two.
Since $a_q,b_q\to0$, in every case $w_q\to0$.
Thus $w\in R_\ell$ and $z=dw$ as actual elements, proving
\[
                         \mathrm{res}([c])=0.
 \tag{A.27.15}
\]
This removes the possible phantom residue. Local divisibility
alone would only have placed $w$ in the full product;
the valuation estimates prove precisely the stronger
membership in $R_\ell$ needed here.

\emph{The unramified term at all odd primes.}
For odd $\ell$, the integral comparison in
\cite[VI, Theorem 8.2]{Weibel2013}, at $n=2$ and
$\mathcal O_S=\Z$, gives
\[
 H^2_{\mathrm{cts}}(G_{\Q,S_\ell},\Z_\ell(2))
   =H^2_{\mathrm{et}}(S_\ell,\Z_\ell(2))
   \simeq K_2(\Z)_{(\ell)}=0.
 \tag{A.27.16}
\]
The last equality uses $K_2(\Z)=\Z/2$, generated by
$\{-1,-1\}$, as in
\cite[III, Example 5.2.2]{Weibel2013}.
The arithmetic comparison and the finite $H^1$ groups
cited above identify continuous cohomology with the
inverse-limit \ensuremath{\ell}-adic cohomology in this
statement. Exactness of (A.27.6) and (A.27.15) give
$A_\ell\cup B_\ell=0$ for every odd $\ell$.

\emph{The integral unramified term at two.}
We give a direct calculation which keeps the real place.
Put $S_2'=\operatorname{Spec}\Z[i,1/2]$.
Its Picard group is zero because $\Z[i]$ is a principal
ideal domain. Its Brauer group is zero: in the global
invariant sequence it can have nonzero invariants only
at the one place above two, and their sum must be zero;
there are no real places. The exact Brauer sequence is
\cite[II, Proposition 2.1]{BIR_MilneADT}, equivalently
\cite[VIII, Theorem 4.2]{MilneCFT} with the unramified
conditions imposed. The Kummer sequence therefore gives
$H^2(S_2',\mathbf F_2)=0$.

The coefficient sequences
\[
0\longrightarrow\Z/2(2)\longrightarrow\Z/2^{n+1}(2)
 \longrightarrow\Z/2^n(2)\longrightarrow0
\]
give $H^2(S_2',\Z/2^n(2))=0$ by induction. The groups
$H^1(S_2',\Z/2^n(2))$ are finite, so their first derived
inverse limit is zero and
\[
 H^2_{\mathrm{cts}}(G_{\Q(i),\{2\}},\Z_2(2))=0.
\]
The corresponding unramified Galois group is the
index-two subgroup of $G_{\Q,S_2}$, since $\Q(i)/\Q$
ramifies only at two. Restriction followed by
corestriction is multiplication by two. Consequently
$H^2_{\mathrm{cts}}(G_{\Q,S_2},\Z_2(2))$ is killed by two.
The exact coefficient sequence for multiplication by two
injects this group into $H^2(S_2,\mathbf F_2)$.
Again the Picard group is zero. The Brauer invariant
sequence identifies
\[
 H^2(S_2,\mathbf F_2)=\operatorname{Br}(S_2)[2]
   =\{(\epsilon_2,\epsilon_\infty):
       \epsilon_2=\epsilon_\infty\in\tfrac12\Z/\Z\}
   \simeq\Z/2.
\]
The upper bound is attained by
$\kappa_2(-1)\cup\kappa_2(-1)$, unramified away from two.
For real complex conjugation $\sigma$, (A.27.13) has value
$-1$. Its cup at $(\sigma,\sigma)$ is $-1$, which represents
the nonzero element of
$H^2(G_\R,\Z_2(2))=\Z_2/2$.
Thus restriction is an isomorphism
\[
 H^2_{\mathrm{cts}}(G_{\Q,S_2},\Z_2(2))
        \xrightarrow{\ \sim\ }H^2(G_\R,\Z_2(2))=\Z/2.
 \tag{A.27.17}
\]

By (A.27.6) and (A.27.15), the global $2$-adic cup comes
from the left side of (A.27.17). Its real restriction is
zero by A.25. One can check this particular fact directly:
its reduction modulo two is the Brauer cup of $A_2,B_2$;
all finite local invariants vanish by the local Steinberg
relation, so the real invariant vanishes by the sum formula.
Reduction
$H^2(G_\R,\Z_2(2))\to H^2(G_\R,\mathbf F_2)$ is an
isomorphism. Therefore the class coming from (A.27.17)
is zero, proving $A_2\cup B_2=0$ as well.

Finally $\widehat\Z(2)=\prod_\ell\Z_\ell(2)$ with its
product topology. A map from any profinite power of $G_\Q$
to this product is continuous exactly when every
coordinate map is continuous. Hence the full continuous
cochain complex is the product of the $\ell$-adic complexes.
Products are exact for abelian groups, so its second
cohomology is the product of their second cohomologies.
The vanishing proved at every $\ell$ gives (A.27.1).
\end{proof}

The calculation supplies the continuous cup itself, with
its first derived inverse-limit component, on all of
$\mathcal G$. A.38 gives the separate identification with
the actual full profinite class-two quotient of the tripod.
A.44 constructs infinite-support pairs in this locus. The
weight-three comparison is given in A.43; the corresponding
vanishing and full metabelian lifting questions remain unresolved. The finite exact checks of the
strict chain identity, the original-cup formula and its
$2$-adic binomial correction are retained in
\texttt{code/cup\_residue\_certificate.py} and
\texttt{code/full\_cup\_class2\_certificate.py}, with their outputs.


\subsection{Triangle stacks, actual maps, and the number-field comparison}

\begin{statement}{A.28}{Proved}{the exact height data supplied by L8}
For $n\ge2$ the root stack
$\mathcal X_n=\Pj^1_\Q[\sqrt[n]{0},\sqrt[n]{1},\sqrt[n]{\infty}]$
has orbifold Euler characteristic $-1+3/n\le1/2$.
It therefore does not satisfy the Euler-characteristic hypothesis
$\chi>1$ of Bhargava--Poonen's Theorem 5. A local tuple, by itself,
does not specify the global point and metrized vector bundle appearing
in Ellenberg--Satriano--Zureick-Brown's Definition 2.11 of height.
In particular the notation $N_n$ in L8 has not been identified there
with such a height.
\end{statement}
\begin{proof}
Each of the three order-$n$ stabilizers subtracts $1-1/n$ from the
Euler characteristic $2$ of the coarse projective line. The stated
formula and inequality follow. The hypothesis of
\cite[Theorem 5]{BhargavaPoonen2020} is genus less than $1/2$,
equivalently Euler characteristic greater than $1$ in this convention.
The displayed calculation is the exact numerical comparison with
that theorem.

At a root-stack chart the coarse coordinate is $t=z^n$ and the chart
is $[\mathbf A^1/\mu_n]$. An ordinary coarse point of valuation
$v(t)$ extends to this chart over $\Z_p$ exactly when the associated
root torsor and line admit an integral presentation; away from primes
dividing $n$, its valuation condition is $n\mid v(t)$. A global
Kummer class modulo $n$ has nonzero valuations modulo $n$ at only
finitely many primes. Thus the fixed-$n$ local points supplied by a
Selmer section extend integrally outside a finite set. At a cuspidal
local section, the same conclusion follows by applying this finite
support assertion to its $\mu_n$-torsor in the residual gerbe.
This proves fixed-level integrality; it supplies no uniform finite set
as $n$ varies.

The height construction of
\cite[Definition 2.11]{ESZB2023} starts with an actual morphism
$\operatorname{Spec}\Q\to\mathcal X_n$ and a metrized vector bundle.
Here the supplied morphisms have sources $\operatorname{Spec}\Q_p$
separately. The missing comparison map would have to produce a
single global point from this family. A product of local norms on an
arbitrary tuple is not an invariant replacement: set $x_2=2$,
$x_p=3$ at every other finite prime, and $x_\infty=3$.
The global unit $t$ on $U$ gives
\[
 \prod_v |t(x_v)|_v=\frac12\cdot\frac13\cdot3=\frac12.
\]
Consequently changing a rational trivialization by $t$ changes that
naive product. For a diagonal rational point the corresponding
factor is $1$ by the product formula. This explicit calculation
locates the missing invariance; it makes no assertion against height
constructions with additional specified data.
\end{proof}

\begin{statement}{A.29}{Proved}{the exact Kummer twist selected at every exponent}
Let $s$ be a full continuous section of $E_U\to G_\Q$, and choose
$a,b\in\Q^\times$ representing $A_s,B_s$ modulo $n$.
Then $s$ lifts to the finite étale cover
\[
 V_{a,b,n}:\quad aX^n+bY^n=1,\quad XY\ne0,
 \qquad t=aX^n.
\]
Its class as a $\mu_n^2$-twist of the specified Kummer cover depends
only on the two classes of $a,b$. The word unique here refers to this
torsor-twist class, with its map to $U$, rather than to a unique
unmarked isomorphism class of projective curves.
\end{statement}
\begin{proof}
The cover obtained by adjoining $n$th roots of $t$ and $1-t$ is a
$\mu_n^2$-torsor over $U$. Its pullback to the section has class
$(A_s,B_s)\in H^1(G_\Q,\mu_n^2)$. Twisting by its negative makes
that pulled-back finite $G_\Q$-set have a fixed point. A fixed point
of the fibre identifies an open subgroup of $E_U$ containing a
representative of $s(G_\Q)$, which is exactly the asserted lift.
The two Kummer equations give the displayed affine model. Replacing
$a,b$ by $a r^n,b u^n$ is the $\Q$-isomorphism
$(X,Y)\mapsto(X/r,Y/u)$ over $U$. Conversely equality of torsor
twists is equality in $H^1(G_\Q,\mu_n^2)$ by the defining descent
cocycles. Forgetting the map and the deck action can identify further
curves; we make no converse for that forgetful operation.
\end{proof}
This supplies the exact finite-level map in L6. The support of the
chosen classes is finite for each $n$. Neither its finiteness nor the
twisting construction gives a bound uniform in $n$. The quadratic
twist correction to the residual assertion R4 is recorded in A.20.

\begin{statement}{A.30}{Proved}{the rational-map action at the birational level}
The monoid of nonconstant rational functions $f\in\Q(t)$ acts, by
pushforward, on continuous birational section classes of
$G_{\Q(t)}\to G_\Q$. The action comes from the field embeddings
$\Q(z)\hookrightarrow\Q(t)$, $z\mapsto f(t)$, and induces the
correspondences of A.8 on model sections. It induces a morphism
$U\to U$ exactly for the six nonconstant maps of A.7.
\end{statement}
\begin{proof}
The field extension $\Q(t)/\Q(f(t))$ is finite separable. In a common
separable closure, restriction gives a continuous open injection
$G_{\Q(t)}\hookrightarrow G_{\Q(f(t))}$ of index $\deg(f)$.
Identify $\Q(f(t))$ with $\Q(z)$ and then rename $z$ as $t$.
This homomorphism respects the quotient $G_\Q$; composition with a
section is a section. Choices of common separable closure change the
map by a geometric inner automorphism, exactly the allowed ambiguity.
Composition of field embeddings gives the claimed action on classes.
For the model, remove from the source all inverse images of
$\{0,1,\infty\}$ and all points at which étaleness is required.
There is then an actual morphism to $U$, and its map on fundamental
groups is the one induced by the field inclusion. The source equals
$U$ only in the nonconstant endomorphism cases classified in A.7.
\end{proof}

\begin{statement}{A.31}{Proved}{the framed fibre over a congruence section}
Let $q:E_U\to E_U/K$ be a continuous quotient with closed normal
$K\subseteq\Pi$, and let $s_0$ be a fixed continuous section.
The set of continuous sections $s$ satisfying $qs=qs_0$, modulo
conjugation by $K$, is canonically the pointed nonabelian cohomology
set $H^1_{\rm cts}(G_\Q,K)$ for the action
$g\cdot k=s_0(g)ks_0(g)^{-1}$. In particular this describes the
framed fibre for the kernel of geometric Legendre monodromy.
\end{statement}
\begin{proof}
For such a section set $c(g)=s(g)s_0(g)^{-1}$. Equality of the
quotient maps places $c(g)$ in $K$. Multiplication of $s(g)s(h)$
gives $c(gh)=c(g)(g\cdot c(h))$. Conversely every continuous
cocycle satisfying this identity defines a continuous section by
$s(g)=c(g)s_0(g)$. Conjugation by $k\in K$ gives exactly
$c(g)\mapsto k c(g)(g\cdot k^{-1})$, the equivalence defining
nonabelian $H^1$. These constructions are inverse.

Forgetting the frame allows the additional conjugations by those
$\gamma\in\Pi$ whose images under $q$ centralize $qs_0(G_\Q)$.
This gives the exact further quotient on the framed fibre. Thus
equality of the monodromy images supplies a cocycle problem in its
specified kernel. It does not compute that cohomology set or replace
it by a family of signs. No such computation is claimed here.
\end{proof}

\begin{statement}{A.32}{Proved}{a finite-support ghost is a rational point}
On $\mathcal G$, finite prime support of either coordinate forces
$(A,B)=(\kappa(r),\kappa(1-r))$ for a unique
$r\in\Q\setminus\{0,1\}$. Its ordinary finite local points are all
equal to $r$.
\end{statement}
\begin{proof}
Suppose $A$ has finite support. Its valuations are ordinary integers
because all its localizations are $\kappa(x_p)$. A.10 gives
$A=\kappa(r)$ for $r\in\Q^\times$. Local Kummer injectivity gives
$x_p=r$ for every finite $p$, and hence $r\ne1$. Consequently all
localizations of $B\kappa(1-r)^{-1}$ are trivial. A.10's valuation
coordinates are therefore all zero. The only remaining coordinate
is its sign, which is detected in $\widehat{\Q_3^\times}$ because
$-1$ is not a square in $\Q_3$. It too is trivial. The argument with
$B$ first is identical. Local injectivity also proves uniqueness.
\end{proof}

\begin{statement}{A.33}{Proved}{why the finite-unit argument over $\Q$ needs modification over a number field}
Let $K=\Q(\sqrt2)$, $\epsilon=1+\sqrt2$, and
$e_3\in\widehat{\Z}=\prod_\ell\Z_\ell$ have $3$-component $1$
and all other components $0$. The element
$A=\epsilon^{e_3}\in\widehat{K^\times}$ is represented by an
ordinary unit of $K_v$ at every finite place $v$, but is not in the
image of $K^\times$. All its finite valuations are zero.
\end{statement}
\begin{proof}
The power map $\Z\to\mathcal O_v^\times$, $m\mapsto\epsilon^m$,
extends continuously to $\widehat{\Z}$ because the local unit group
is profinite. Its value at $e_3$ represents the localization of $A$.
This proves the local assertion, including at places above $2$ and
$3$.

For completeness $\mathcal O_K=\Z[\sqrt2]$ and
$\mathcal O_K^\times=\{\pm\epsilon^m:m\in\Z\}$. The latter can
be checked directly. Multiply a positive real unit by a suitable
power of $\epsilon$ so that $1\le u<\epsilon$. Its conjugate is
$\pm u^{-1}$. For norm $+1$, its $\sqrt2$ coefficient
$(u-u^{-1})/(2\sqrt2)$ is an integer in $[0,1)$, so is zero.
For norm $-1$, that coefficient is
$(u+u^{-1})/(2\sqrt2)$, strictly between $0$ and $1$, which is
impossible. Thus the only reduced unit is $1$.
The divisor exact sequence embeds the unit group as a direct summand
of $K^\times$: its quotient is the group of principal divisors,
a subgroup of a free abelian group and hence free abelian. Exponent
completion consequently preserves this embedding.
If $A$ were represented by $r\in K^\times$, its zero valuations
would give $r=\pm\epsilon^m$. In the completed free unit factor we
would then have $m=e_3$. Its $5$-adic component gives $m=0$ as an
ordinary integer, whereas its $3$-adic component gives $m=1$.
This contradiction proves the assertion.
\end{proof}
The example gives an explicit comparison for A11: the implication
from local ordinary units and finite valuation support to a global
number used in A.32 is valid over $\Q$; over this $K$ its completed
unit factor already contains the displayed additional element.
It is not a pair satisfying the completed unit equation and is not
a counterexample section.
